% -----------------------------------------------------------------------------
% Beginner bridge: how ordinary one-particle quantum mechanics emerges
% -----------------------------------------------------------------------------

量子场算符可以改变粒子数，而普通量子力学的波函数只是复数值函数；两种语言之间的联系最好先在一个守粒子数的非相对论 Schr\"odinger 场中看清。这个低能模型把 Fock 空间按粒子数分解，并使普通单粒子量子力学直接作为其中的固定一粒子扇区出现。

等时正则关系为
\begin{equation}
 [\hat\Psi(\vb{r}),\hat\Psi^\dagger(\vb{r}')]
 =\delta^{(3)}(\vb{r}-\vb{r}'),
 \qquad
 [\hat\Psi,\hat\Psi]=0.
 \label{eq:sch-field-comm-r6}
\end{equation}
若真空满足 $\hat\Psi(\vb{r})|0\rangle=0$，则 $\hat\Psi^\dagger(\vb{r})|0\rangle$ 是在位置标签 $\vb{r}$ 上产生一个粒子的广义态。任意一粒子态因此可展开为
\begin{equation}
 {
 |\Psi_1(t)\rangle
 =\int\dd^3 r\,
 \psi_{\rm 1p}(\vb{r},t)
 \hat\Psi^\dagger(\vb{r})|0\rangle .}
 \label{eq:fock-oneparticle-expansion-dp68}
\end{equation}
这里的 $\psi_{\rm 1p}$ 才是普通第一量子化中的复数值波函数；$\hat\Psi$ 本身仍是作用在 Fock 空间上的算符值分布。

对自由哈密顿量
\begin{equation*}
 \hat H
 =\int\dd^3 r\,
 \hat\Psi^\dagger(\vb{r})
 \left(-\frac{\hbar^2\nabla^2}{2m}\right)
 \hat\Psi(\vb{r}),
\end{equation*}
Schr\"odinger 绘景的态方程在一粒子扇区中约化为


\begin{proposition}[二次量子化在一粒子空间上的限制]
设一粒子算子为自由非相对论哈密顿量 $h=-\hbar^2\nabla^2/(2m)$，取本节指定边界条件下的自伴定义域，并令 $\hat H=\dd\Gamma(h)$。一粒子态及其作用均属于相应定义域，下式成立；一般自伴 $h$ 的对应一粒子方程为 $\ii\hbar\partial_t\psi=h\psi$：
\begin{equation}
 {
 \ii\hbar\partial_t\psi_{\rm 1p}(\vb{r},t)
 =-\frac{\hbar^2}{2m}\nabla^2\psi_{\rm 1p}(\vb{r},t).}
\label{eq:fock-oneparticle-schrodinger-dp68}
\end{equation}
\end{proposition}
\begin{proof}
一粒子波函数不是额外附加到场论上的对象，而是 Fock 态在局域一粒子基底上的坐标。先从正文的态展开
\begin{equation*}
|\Psi_1(t)\rangle
=\int\dd^3r\,\psi_{\rm 1p}(\vb{r},t)
\hat\Psi^\dagger(\vb{r})|0\rangle
\end{equation*}
出发。自由 Schr\"odinger 场哈密顿量为
\begin{equation*}
\hat H
=\int\dd^3x\,
\hat\Psi^\dagger(\vb{x})
\left(-\frac{\hbar^2\nabla_{\vb{x}}^2}{2m}\right)
\hat\Psi(\vb{x}).
\end{equation*}
利用等时对易关系，先求产生一个局域粒子时能量算符怎样作用：
\begin{align*}
[\hat H,\hat\Psi^\dagger(\vb{r})]
&=\int\dd^3x\,
\hat\Psi^\dagger(\vb{x})
\left(-\frac{\hbar^2\nabla_{\vb{x}}^2}{2m}\right)
[\hat\Psi(\vb{x}),\hat\Psi^\dagger(\vb{r})]\\
&=\int\dd^3x\,
\hat\Psi^\dagger(\vb{x})
\left(-\frac{\hbar^2\nabla_{\vb{x}}^2}{2m}\right)
\delta^{(3)}(\vb{x}-\vb{r})\\
&=-\frac{\hbar^2}{2m}\nabla_{\vb{r}}^2
\hat\Psi^\dagger(\vb{r}).
\end{align*}
因为真空满足 \(\hat H|0\rangle=0\)，于是
\begin{align*}
\hat H|\Psi_1\rangle
&=\int\dd^3r\,\psi_{\rm 1p}(\vb{r})
[\hat H,\hat\Psi^\dagger(\vb{r})]|0\rangle\\
&=\int\dd^3r\,
\left[-\frac{\hbar^2}{2m}\nabla^2\psi_{\rm 1p}(\vb{r})\right]
\hat\Psi^\dagger(\vb{r})|0\rangle,
\end{align*}
其中第二行对空间变量分部积分两次，并假设波函数在边界处足够快衰减。

连续基底中的系数比较由场算符投影严格实现。对任意 \(\vb{x}\)，左乘 \(\langle0|\hat\Psi(\vb{x})\)，由
\begin{equation*}
\langle0|\hat\Psi(\vb{x})
\hat\Psi^\dagger(\vb{r})|0\rangle
=\delta^{(3)}(\vb{x}-\vb{r})
\end{equation*}
可得
\begin{align*}
\langle0|\hat\Psi(\vb{x})
\ii\hbar\partial_t|\Psi_1\rangle
&=\ii\hbar\partial_t\psi_{\rm 1p}(\vb{x},t),\\
\langle0|\hat\Psi(\vb{x})
\hat H|\Psi_1\rangle
&=-\frac{\hbar^2}{2m}\nabla_{\vb{x}}^2
\psi_{\rm 1p}(\vb{x},t).
\end{align*}
因此 Fock 空间 Schr\"odinger 方程 \(\ii\hbar\partial_t|\Psi_1\rangle=\hat H|\Psi_1\rangle\) 在每个 \(\vb{x}\) 上严格投影成
\begin{equation*}
\ii\hbar\partial_t\psi_{\rm 1p}(\vb{x},t)
=-\frac{\hbar^2}{2m}\nabla_{\vb{x}}^2\psi_{\rm 1p}(\vb{x},t),
\end{equation*}
即正文\eqnrefs{eq:fock-oneparticle-schrodinger-dp68}。这里真正建立的是“固定粒子数 Fock 扇区”与普通单粒子波函数之间的坐标同构，而不是重新假设一条量子力学方程。
\end{proof}

\eqnrefs{eq:fock-oneparticle-schrodinger-dp68} 不是另加的一条动力学假设，而是把场论哈密顿量限制到\eqnrefs{eq:fock-oneparticle-expansion-dp68} 所定义的一粒子扇区后的结果。

$\hat\Psi^\dagger$ 用来构造粒子态，$\psi_{\rm 1p}$ 用来描述固定一粒子态在位置基中的展开系数。进入两粒子及更高扇区后，态由多个产生算符构造，玻色对易关系或费米反对易关系自动决定交换对称性，因此不需要再把场算符误解成一条“多粒子波函数”。


固定粒子数扇区解释了普通波函数从哪里来。下一步回到相对论场论：此时粒子数不再必须守恒，真正需要保留的是局域场及其二点函数；类空分离处的场对易关系将给出微观因果性。

